% page 554 du fac-simile (image p-553 du scan)
% bloc 157.5 x 242.0 mm ; marge gauche 8.0 mm
\pgc{-12.700mm}{0.000mm}{
\l{\cel{3}546}
\l{}
\l{\sou{spico}. Singla ek la substanci vejetastala, aromata o pikan-}
\l{\cel{3}ta, quin onu uzas por saporozigar disho : kariofilo, mus-}
\l{\cel{3}kado, pipro, cinamo, jinjero pulverigita. - DEFIS.}
\l{}
\l{\sou{spiko}. (\sou{bot.)}\cel{2}Parto extrema di la stipo di la graminei, qua}
\l{\cel{3}portas, asemblita cirkum axo, la grani di la planto.- EIS}
\l{}
\l{spino. (\sou{anat.}) (\sou{che la homo ed ula animali altra)}. Ensemblo}
\l{\cel{3}ek la vertebri dorsala, artikoza, ye qui konsistas la axo}
\l{\cel{3}di la karpenturo osta. - EFIS.}
\l{}
\l{\sou{spinato}. (\sou{bot.)}\cel{2}Planto ek la familio \textquotedbl{}kenopodiacei\textquotedbl{}, herbo}
\l{\cel{3}un-yar-dura, di qua la folii, pos koqueso, esas nutrivo}
\l{\cel{3}bonega. - DEIRS.}
\l{}
\l{\sou{spindelo}. I.Stango ligna, di qua la extremaji esas dinigita,}
\l{\cel{3}un de li pinta, la altra rondigita, quan onu uzas lor fil-}
\l{\cel{3}ifo per ruko, por tordar la filo e rular lu segun ke lu}
\l{\cel{3}formacesas. - II. (\sou{roto di vehilo}). Parto di la rot-axo}
\l{\cel{3}qua fricionas kontre la kuseneti. - DE.}
\l{}
\l{\sou{spinelo}. (\sou{mineral.}) Rubino pale-reda. -}
\l{}
\l{\sou{spionar}.\cel{2}(\sou{trans.}) Observar habile e sekrete ta (o ti) di}
\l{\cel{3}qua (o di qui) onu profitos de saveskar lua (o lia) inten-}
\l{\cel{3}ci - o to quo eventas che la enemiko. - DEFIRS.}
\l{}
\l{\sou{spiro}. Lineo qua volvas cirkum axo quale la fileto di skrubo,}
\l{\cel{3}di helico. - DEFIRS.}
\l{}
\l{\sou{spiralo}. I. (\sou{geom.}) Kurvolineo, sur plano, qua trasas sama-}
\l{\cel{3}sinse serio de rotaci cirkum punti fixa (polo) de qua lu}
\l{\cel{3}deviacas gradoze segun +leyo (lego) determinita. - II.}
\l{\cel{3}(\sou{horloj.}) Mikra resorto qua reglas la eskapo di la pendu-}
\l{\cel{3}lo de horlojeto, e mantenas la izokroneso di lua ocili. -}
\l{\cel{3}DEFIRS.}
\l{}
\l{\sou{spireo}. (\sou{bot.})\cel{2}Genero de planti ek la familio \textquotedbl{}rozacei\textquotedbl{}.-}
\l{\cel{3}L.\cel{1}\sou{spiraeum}.- DEFIS.}
\l{}
\l{spiremo. (\sou{biol}.)Reto de nuklei qua, kande mitoso esas quik}
\l{\cel{3}eventonta, divenas, komence, kordono longa, tenua, dornoza,}
\l{\cel{3}kun falduri tre interproxima - e, pose, divenas plu kurta,}
\l{\cel{3}plu dika, kun surfaco glata, e di qua la falduri min-}
\l{\cel{3}proximeskas. - DEFIS.}
\l{}
\l{\sou{spirito}. (\sou{religio)}. Substanco nekorpala (kontraste kun ma-}
\l{\cel{3}\sou{terio}, substanco korpala) a qua uli imputas la penso-fa-}
\l{\cel{3}kultato. - DEFIRS.}
\l{}
\l{\sou{spiso}. Stango ek fero, di qua un extremajo esas pinta, quan}
\l{\cel{3}onu pasigas tra grosa peco de karno, tra pultruno, por}
\l{\cel{3}rostar li, e quan onu rotacigas, o manue, o per mekan-}
\l{\cel{3}ismo. - DES.}
}
